\section{Experimental Protocol}
\label{sec:experiments}

Each of the 13 controllers runs on the physical platform under the same conditions: the sensors, actuator, and 20\,Hz control loop detailed in Sec.~\ref{sec:system}, under the stratified initial-condition schedule below.

\subsection{Stratified Initial Conditions}

To probe different failure modes, we stratify the 50 trials by difficulty.
Difficulty is set mainly by the cart's starting position, and the two at-wall strata split further by which side of the cart the ball starts on.

\noindent\textbf{Center} (10 trials): cart near center ($\pm$0.10\,m), maximum maneuverability in both directions.
\noindent\textbf{Mid-rail} (16 trials): cart at $\pm$0.40\,m, with asymmetric maneuvering space.
\noindent\textbf{Near-wall} (14 trials): cart at $\pm$0.65\,m, leaving only 0.13\,m of clearance to the boundary.
\noindent\textbf{At-wall-opposite} (5 trials): cart at the boundary ($\pm$0.78\,m), ball displaced to the side away from the wall.
Getting under the ball drives the cart back toward open rail, so the correction has room to run; the easier of the two at-wall cases.
In this stratum the cart did not always reach the physical limit, so the achieved start was not uniform across controllers (supplement, Evaluation Protocol); we therefore read it as a recovery-speed check rather than an equal-condition comparison.
\noindent\textbf{At-wall-same} (5 trials): cart at the boundary ($\pm$0.78\,m), ball displaced to the same side as the cart.
Getting under the ball means driving further into the wall, so this stratum stresses controllers that lack explicit constraint handling; the harder of the two at-wall cases.


\subsection{Success Criterion}

A trial succeeds when the ball angle stays within $\pm$0.01\,rad continuously for 1.0\,s; a trial that runs past 30\,s is a failure.
At the arc radius $R = 2.101$\,m of Sec.~\ref{sec:system}, $\pm$0.01\,rad is about $\pm$21\,mm of arc, against sensor readings that run about 13\,mm peak-to-peak.

\subsection{Metrics}

We report success rate (SR) with 95\% Wilson score intervals, median settling time with interquartile range over the successful trials, and RMS control effort (the root-mean-square commanded cart velocity in m/s, saturated at the $\pm 0.9$~m/s axis limit that the actuator enforces, then averaged over all trials).
For pairwise comparisons we use the Mann-Whitney rank test on the successful-trial settling times, with Bonferroni correction across all 78 controller pairs; a rank test suits these skewed distributions better than a mean-based $t$-test.
Success-rate differences are compared with Fisher's exact test in the same Bonferroni family.
We use 50 trials per controller as a balance between confidence-interval width and hardware-time budget: each trial runs on the physical rig, up to 30\,s of balancing plus an automated reset that jolts the cart to reposition the ball for the next trial.
Settling statistics are computed over successful trials only, so cross-controller settling comparisons slightly favor the less-reliable controllers, whose failed (and typically hardest) trials are excluded; success rate is therefore reported separately and never folded into the settling comparison.

\subsection{A Note on Statistical Power}

With 50 trials per controller, the benchmark resolves only large differences in success rate, and its pairwise tests hold the family-wise error rate across all 78 controller pairs.
At that corrected level, a Fisher comparison reaches 80\% power only for success-rate gaps of roughly 30 percentage points near the top of the range, widening to about 40 points near 50\%; for a single uncorrected comparison the figures are closer to 20 and 30 points.
Gaps of a few points sit below this resolution.

The confidence intervals show the same limit directly.
A controller that succeeds in all 50 of its trials still carries a 95\% Wilson interval of $[93, 100]$, and one that succeeds in 40 carries $[67, 89]$, so any single success rate is uncertain by several points on each side.

We chose 50 trials rather than the more common 10 or 20 because the interval narrows substantially with sample size: at ten trials the interval on a 90\% success rate is about twice as wide as at fifty.
Collecting the full 650-trial dataset (13 controllers $\times$ 50 trials) took several sessions at the rig, on top of thousands of development and screening trials.

\subsection{Auxiliary Studies}

A few auxiliary studies use the same platform and protocol and are reported where they inform a finding: the wall-override ablations (each classical controller with and without the boundary rule), the PPO domain-randomization and blind-zone ablations, an open-loop sweep of world-model accuracy against dataset size, and a short on-hardware fine-tuning trial.
The supplement gives their full protocols.
